\chapter{Introducción}

En las últimas décadas se ha presenciado un aumento exponencial en la creación e intercambio de información por los usuarios en plataformas digitales, fenómeno conocido como \textit{electronic Word of Mouth} (eWOM). Esta dinámica ha transformado la manera en que las organizaciones comprenden y responden a las necesidades y preferencias de sus clientes respecto a un producto o servicio concreto \citep{mishra2016ewom}. En la industria hotelera, se han realizado numerosos estudios que evidencian el impacto significativo del eWOM en la percepción de los consumidores y su influencia en la toma de decisiones (tanto de los potenciales usuarios como de los propios proveedores de servicios) \citep{cantallops2014new}. Un comportamiento similar se observa en el sector restaurantero, donde las opiniones y reseñas son cruciales para conocer la percepción de los clientes sobre la calidad de los alimentos, el servicio y la experiencia general con el fin de desarrollar estrategias de administración efectivas \citep{bilgihan2018identifying, ong2012perceived}.



\section{Hipótesis de trabajo} \label{sec:hypothesis}

\section{Objetivos} \label{sec:objective}

\section{Contribución} \label{sec:contribution}

En literatura existe una basta colección de trabajos que abordan el problema...

\section{Organización de la tesis}






